\chapter{Giới thiệu}
\label{sec:p1-gioi-thieu}
% Bản nháp P1: mục tiêu/phạm vi, chưa khẳng định kết quả thực nghiệm.
Kiểm chứng thiết kế (verification) đánh giá thiết kế có đáp ứng đặc tả hay không,
trong khi kiểm thử sau sản xuất (manufacturing testing) nhằm phát hiện các chip
có khuyết tật chế tạo. Một thiết kế đúng vẫn có thể tạo ra sản phẩm lỗi, do đó
hai hoạt động này bổ sung cho nhau. Thay vì kiểm tra mọi tổ hợp đầu vào, kiểm
thử cấu trúc sử dụng mô hình lỗi để xây dựng các mẫu kiểm tra có mục tiêu.
Sinh mẫu kiểm tra tự động (Automatic Test Pattern Generation -- ATPG) tìm
các kích thích làm cho đáp ứng của mạch lỗi khác mạch tốt tại ngõ ra
quan sát được~\cite{p1_wang2006}.

Đề tài tập trung giải thích nguyên lý ATPG, cơ chế của D-algorithm và PODEM,
đồng thời trình bày những khó khăn khi áp dụng ATPG cho mạch tuần tự.
Phần thực hành hướng tới cài đặt PODEM bằng Python, minh họa trên mạch c17
và đối chiếu mẫu sinh ra bằng mô phỏng lỗi. Lỗi \texttt{11/SA0} được chọn
làm ví dụ chung để so sánh quá trình tìm kiếm bằng tay với chương trình.

Phạm vi chính là mô hình lỗi kẹt mức logic đơn (single stuck-at fault).
Đối với mạch tuần tự nhỏ, nhóm khảo sát thiết kế quét đầy đủ (full scan)
và trải khung thời gian (time-frame expansion). Các phần tiếp theo lần lượt
trình bày cơ sở lý thuyết, thuật toán, cách cài đặt, kết quả kiểm chứng và
kết luận; những giới hạn thực nghiệm sẽ được nêu theo kết quả thực tế.
